\documentclass[10pt]{article}
\usepackage[margin=0.6in]{geometry}
\usepackage{xcolor}
\usepackage{tcolorbox}
\usepackage{hyperref}
\usepackage{enumitem}
\usepackage{booktabs}

% Define colors
\definecolor{osaorange}{RGB}{220,115,38}
\definecolor{aiblue}{RGB}{30,144,255}

% Configure hyperref
\hypersetup{
    colorlinks=true,
    linkcolor=blue,
    urlcolor=blue,
    citecolor=blue
}

% Compact spacing
\setlength{\parindent}{0pt}
\setlength{\parskip}{4pt}

\title{\textbf{\textcolor{osaorange}{H524 Group Project Check-In: AI Divergence Experiment}}\\
Fall 2025}
\author{Dr. John Molitor | Due: See Canvas | 5\% of Project Grade}
\date{}

\begin{document}
\maketitle
\vspace{-2em}

\begin{tcolorbox}[colback=aiblue!10,colframe=aiblue,boxrule=1.5pt,arc=2pt,left=4pt,right=4pt,top=3pt,bottom=3pt]
\small
\textbf{Purpose:} Discover how AI tools respond differently to biostatistical questions and learn to use prompt engineering to reach consensus.
\end{tcolorbox}
\vspace{-0.5em}

\section*{The Experiment (4 Steps)}

\textbf{Step 1: Propose a Research Question} - Choose a \textbf{broad question} about your dataset that might elicit different AI answers. Good questions ask about \textit{method choice, assumptions, or interpretation}. Examples: "What statistical method for treatment effects in pre-post data?" or "How to handle three-group comparisons?"

\textbf{Step 2: Consult Two Different AI Tools} - Use the \textbf{EXACT same prompt} with two AI tools (Claude, ChatGPT, Copilot, or Gemini). Copy their full responses.

\textbf{Step 3: Use Third AI as Judge} - Ask a different AI: \textit{"I asked two AI tools [your question] and got these responses: [paste both]. Rate their similarity 0-100 (0=different, 100=identical). Identify key agreements/disagreements."}

\textbf{Step 4: Iterate to Increase Consensus} - Try $\geq$2 modified prompts (add context, be specific, add constraints, request format). Re-score similarity each time.

\section*{Deliverable: 1-2 Page Report}

\textbf{Section 1:} Team info (number, members, dataset)

\textbf{Section 2:} Your research question and why you expected divergence

\textbf{Section 3:} Experiment table showing $\geq$3 iterations with: Prompt text, AI 1 response summary, AI 2 response summary, Similarity score (0-100)

\textbf{Section 4:} Brief reflection (4-5 sentences): Did similarity increase? What prompt changes worked? What did you learn about AI variability? Which AI will you use for final project?

\section*{Grading (5 pts = 5\% of Project)}

\begin{small}
\begin{tabular}{lc}
\toprule
\textbf{Component} & \textbf{Pts} \\
\midrule
Question broad enough for divergence & 1 \\
Consulted 2 AIs with identical prompt & 1 \\
Used 3rd AI to score similarity & 1 \\
$\geq$2 prompt iterations documented & 1 \\
Thoughtful reflection & 1 \\
\midrule
\textbf{Total} & \textbf{5} \\
\bottomrule
\end{tabular}
\end{small}

\vspace{0.3em}
\begin{tcolorbox}[colback=aiblue!5,colframe=aiblue,boxrule=1pt,arc=2pt,left=4pt,right=4pt,top=3pt,bottom=3pt]
\small
\textbf{Why This Matters:} Your final project (30\% of grade) requires comparing AI tools critically. This teaches you that AI responses vary by tool AND prompt. Mastering prompt engineering now strengthens your analysis later.
\end{tcolorbox}

\end{document}
